\documentclass{article}
\usepackage{booktabs}
\usepackage{tabularx}
\usepackage[margin=1in]{geometry}
\begin{document}

\begin{table}[tbp] \centering
\newcolumntype{R}{>{\raggedleft\arraybackslash}X}
\newcolumntype{L}{>{\raggedright\arraybackslash}X}
\newcolumntype{C}{>{\centering\arraybackslash}X}

\caption{Formal hypothesis tests: H3, H6, NE-adjacency (C.9), and horticulture co-explanatory (C.7)}
\label{tab:formal_hypotheses}
\begin{tabularx}{\linewidth}{@{}lCCCCCC@{}}

\toprule
&{(1)}&{(2)}&{(3)}&{(4)}&{(5)}&{(6)} \tabularnewline \midrule
{Variable}&{}&{}&{}&{}&{}&{} \tabularnewline
\midrule \addlinespace[\belowrulesep]
adj\_neuchatel\_coef&&&&&7.71&
adj\_neuchatel\_stderr&&&&&(7.90)&
parcel\_c\_coef&&&2.97&1.68&&
parcel\_c\_stderr&&&(15.1)&(19.5)&&
protestant\_c\_coef&&--548&&--572&&
protestant\_c\_stderr&&(714)&&(1,024)&&
vineyard\_X\_parcel\_coef&&&394&282&&
vineyard\_X\_parcel\_stderr&&&(2,407)&(2,974)&&
vineyard\_X\_protestant\_coef&&359&&325&&
vineyard\_X\_protestant\_stderr&&(483)&&(888)&&
vineyard\_c\_coef&484**&472**&1,096&873&557***&478**
vineyard\_c\_stderr&(200)&(206)&(3,657)&(4,597)&(176)&(202)
horticulture\_per\_cap\_coef&&&&&&6.32
horticulture\_per\_cap\_stderr&&&&&&(7.99)
french\_share\_coef&--28.3***&--26.4***&--31.4&--28.5&--32.7***&--28.2***
french\_share\_stderr&(7.11)&(7.75)&(26.8)&(34.4)&(5.48)&(8.64)
catholic\_share\_coef&8.43*&--539&8.03&--563&9.86*&11.5*
catholic\_share\_stderr&(4.59)&(714)&(5.55)&(1,022)&(4.86)&(5.99)
\_cons\_coef&64.6***&376&68.1***&392&63.6***&58.4***
\_cons\_stderr&(3.35)&(406)&(20.4)&(574)&(3.66)&(8.73)
N&25&25&25&25&25&25
r2&0.538&0.564&0.546&0.565&0.572&0.558
spec&H\_KEY\_centered&H3\_coalition\_strat&H6\_olsonian\_interaction&H3H6\_joint\_strat&H\_KEY\_adj\_NE&H\_KEY\_horticulture
model&ols&ols&ols&ols&ols&ols
dropped&&&&&&
\bottomrule \addlinespace[\belowrulesep]

\end{tabularx}
\\ \parbox{\linewidth}{\footnotesize Notes: Formal tests of two implications of the bootleggers-and-baptists framework derived from Becker (1983) and Olson (1965), pre-specified design. Column (1) replicates the headline KEY spec with a mean-centered vineyard regressor (coefficient identical to the uncentered version). Column (2) tests H3 (coalition interaction): wine-industry effect amplified by moralist coalition strength, operationalized as protestant\\_share\\_total = 1 - catholic\\_share\\_total in the absence of canton-level Blue Cross / Croix-Bleue / IOGT membership data. Column (3) tests H6 (Olsonian concentration): wine-industry effect amplified by industrial concentration, operationalized as avg\\_parcel\\_area\\_1905 (constructed from agricultural land 1912 / total parcels 1905; documented in CONTEXT.md). Column (4) tests both interactions jointly. Column (5) is the C.9 NE-adjacency robustness: KEY spec + adj\\_neuchatel (1 if canton borders Neuchatel: BE, VD, FR). Tests whether the headline vineyard effect survives controlling for spatial spillover from the absinthe-producing canton; column (5) shows it amplifies (vine\\_c +557) rather than attenuates, consistent with NE-adjacency being a wine-industry signal rather than an absinthe-cluster confound. Column (6) is the C.7 multi-Bootlegger horticulture co-explanatory: KEY spec + horticulture\\_per\\_cap (Gartenbau enterprises per 1000 pop, 1905). \\textbf\{Result is INCONCLUSIVE / backmatter\}: horticulture coefficient is positive but statistically null (\\beta = +6.32, SE = 7.99). The Gartenbau measure aggregates vegetable gardens, ornamentals, fruit orchards, and plant nurseries; only the fruit-orchard component directly produces Obstler substrate. The null is consistent with both (a) absence of a fruit-spirit Bootlegger coalition member and (b) measurement noise from the heterogeneous Gartenbau category. Formal multi-Bootlegger test deferred to companion paper P1B with firm-level data. Caveat on the pre-specified design: columns (2) and (4) include BOTH catholic\\_share and protestant\\_c as religion controls. In 1900 Switzerland Catholic + Protestant constituted 99.4 percent of total population, so these two religion variables are near-mechanically collinear (main-effect VIFs above 25\{,\}000 in the joint specification). The interaction coefficient itself is identified despite this main-effect collinearity (interactions and their components do not share variance the same way), so the H3 result reported here remains substantively interpretable. The religion main effects in columns (2) and (4) should not be interpreted. A cleaner variant dropping catholic\\_share from columns (2) and (4) is reported in the backmatter Table\~\\ref\{tab:formal\_hypotheses\_alt\} for transparency. Centered regressors (vineyard\\_c, protestant\\_c, parcel\\_c) ease interpretation: the main effect of vineyard\\_c is the vineyard slope at the mean of the moderator(s) in that spec; the interaction coefficient is the change in that slope per unit increase in the (mean-centered) moderator. HC3 robust SEs in parentheses. N=25 cantons. * p<0.10, ** p<0.05, *** p<0.01.}
\end{table}
\end{document}
